\begin{question}[index=17,points=10]
设可微函数 $f(x,y)$ 满足 $\mathrm df(x,y)=(3x^2-2ax)\dd x+(3y^2-2ay)\dd y$（$a>0$），且 $f(0,0)=0$.
\par（Ⅰ）求 $f(x,y)$；
\par（Ⅱ）若 $f(x,y)$ 有极小值 $-8$，求 $a$ 的值.
\end{question}
\QuestionBreak\WorkSpace

\begin{question}[index=18,points=12]
设 $f(x)$ 满足 $(1-x)f'(x)+nf(x)=n(1-x)^{n+1}$（$n$ 为正整数），且 $f(0)=0$.
\par（Ⅰ）若 $f(x)$ 在 $(0,1)$ 内的最大值为 $M(n)$，求 $\displaystyle\lim_{n\to\infty}M(n)$；
\par（Ⅱ）在第（Ⅰ）问的基础上，证明：$M(n)<\dfrac1\eu$.
\end{question}
\QuestionBreak\WorkSpace\WorkPageBreak

\begin{question}[index=19,points=12]
设 $f(x)$ 有连续的导数，$L$ 为从点 $O(0,0)$ 沿 $x^2+(y-\pi)^2=\pi^2$ 右半圆周到点 $A(0,2\pi)$ 的一段弧，计算
\[I=\int_L\frac{f'(x)\sin y\dd x+[f(x)\cos y-\pi x]\dd y}{x^2+(y-\pi)^2}.\]
\end{question}
\QuestionBreak\WorkSpace

\begin{question}[index=20,points=12]
设 $f(x)$ 在 $(-\infty,+\infty)$ 内有二阶连续导数，且对任意 $x,y\in(-\infty,+\infty)$，有 $f\!\left(\dfrac{x+y}2\right)\le\dfrac{f(x)+f(y)}2$.
\par（Ⅰ）证明：对任意 $x\in(-\infty,+\infty)$，有 $f''(x)\ge0$；
\par（Ⅱ）若 $f(x)$ 在 $[-1,1]$ 上为非负偶函数，证明：$\displaystyle\int_{-1}^1f(x)\dd x\le f(0)+f(1)$.
\end{question}
\QuestionBreak\WorkSpace\WorkPageBreak

\begin{question}[index=21,points=12]
设矩阵 $\vect A=\mat{1&a&a\\a&1&a\\a&a&1}$ 与 $\vect B=\mat{1&1&0\\1&1&0\\0&0&3}$ 合同.
\par（Ⅰ）求 $a$ 的值；
\par（Ⅱ）求一个可逆矩阵 $\vect C$，使得 $\vect C^{\mathrm T}\vect A\vect C=\vect B$；
\par（Ⅲ）在第（Ⅱ）问的基础上，若 $\vect X=(x_1,x_2,x_3)^{\mathrm T}$ 满足 $\operatorname{tr}[\vect A(\vect C\vect X)(\vect C\vect X)^{\mathrm T}]=0$，求 $\vect X$.
\end{question}
\QuestionBreak\WorkSpace

\begin{question}[index=22,points=12]
设二维随机变量 $(X,Y)$ 的概率密度为 $f(x,y)=\begin{cases}kx\eu^{-x(1+y)},&x>0,y>0,\\0,&\text{其他}.\end{cases}$
\par（Ⅰ）求 $k$ 及 $E(X)$；
\par（Ⅱ）求 $P\{XY>2\mid XY>E(XY)\}$.
\end{question}
\WorkSpace
